% Options for packages loaded elsewhere
\PassOptionsToPackage{unicode}{hyperref}
\PassOptionsToPackage{hyphens}{url}
\documentclass[
  polish,
  11pt,
  a4paper,
]{article}
\usepackage{xcolor}
\usepackage[margin=2cm]{geometry}
\usepackage{amsmath,amssymb}
\setcounter{secnumdepth}{-\maxdimen} % remove section numbering
\usepackage{iftex}
\ifPDFTeX
  \usepackage[T1]{fontenc}
  \usepackage[utf8]{inputenc}
  \usepackage{textcomp} % provide euro and other symbols
\else % if luatex or xetex
  \usepackage{unicode-math} % this also loads fontspec
  \defaultfontfeatures{Scale=MatchLowercase}
  \defaultfontfeatures[\rmfamily]{Ligatures=TeX,Scale=1}
\fi
\usepackage{lmodern}
\ifPDFTeX\else
  % xetex/luatex font selection
\fi
% Use upquote if available, for straight quotes in verbatim environments
\IfFileExists{upquote.sty}{\usepackage{upquote}}{}
\IfFileExists{microtype.sty}{% use microtype if available
  \usepackage[]{microtype}
  \UseMicrotypeSet[protrusion]{basicmath} % disable protrusion for tt fonts
}{}
\makeatletter
\@ifundefined{KOMAClassName}{% if non-KOMA class
  \IfFileExists{parskip.sty}{%
    \usepackage{parskip}
  }{% else
    \setlength{\parindent}{0pt}
    \setlength{\parskip}{6pt plus 2pt minus 1pt}}
}{% if KOMA class
  \KOMAoptions{parskip=half}}
\makeatother
\usepackage{longtable,booktabs,array}
\usepackage{caption}
\captionsetup[table]{skip=6pt}
\newcounter{none} % for unnumbered tables
\usepackage{calc} % for calculating minipage widths
% Correct order of tables after \paragraph or \subparagraph
\usepackage{etoolbox}
\makeatletter
\patchcmd\longtable{\par}{\if@noskipsec\mbox{}\fi\par}{}{}
\makeatother
% Allow footnotes in longtable head/foot
\IfFileExists{footnotehyper.sty}{\usepackage{footnotehyper}}{\usepackage{footnote}}
\makesavenoteenv{longtable}
\usepackage{graphicx}
\makeatletter
\newsavebox\pandoc@box
\newcommand*\pandocbounded[1]{% scales image to fit in text height/width
  \sbox\pandoc@box{#1}%
  \Gscale@div\@tempa{\textheight}{\dimexpr\ht\pandoc@box+\dp\pandoc@box\relax}%
  \Gscale@div\@tempb{\linewidth}{\wd\pandoc@box}%
  \ifdim\@tempb\p@<\@tempa\p@\let\@tempa\@tempb\fi% select the smaller of both
  \ifdim\@tempa\p@<\p@\scalebox{\@tempa}{\usebox\pandoc@box}%
  \else\usebox{\pandoc@box}%
  \fi%
}
% Set default figure placement to htbp
\def\fps@figure{htbp}
\makeatother
\ifLuaTeX
\usepackage[bidi=basic,shorthands=off]{babel}
\else
\usepackage[bidi=default,shorthands=off]{babel}
\fi
\ifLuaTeX
  \usepackage{selnolig} % disable illegal ligatures
\fi
\setlength{\emergencystretch}{3em} % prevent overfull lines
\providecommand{\tightlist}{%
  \setlength{\itemsep}{0pt}\setlength{\parskip}{0pt}}
\usepackage{fontspec}
\defaultfontfeatures{Ligatures=TeX}
\usepackage{amsmath,amssymb}
\usepackage{booktabs}
\usepackage{array}
\usepackage{geometry}
\usepackage{enumitem}
\usepackage{xcolor}
\usepackage{tcolorbox}
\usepackage{fancyhdr}
\usepackage{titlesec}
\usepackage{multicol}

\geometry{margin=1.8cm, top=2.5cm, bottom=2cm}
\pagestyle{fancy}
\fancyhf{}
\rhead{\small Projektowanie badań — 2026/27}
\lhead{\small Zarządzanie informacją | ISI UJ}
\cfoot{\thepage}

\titleformat{\section}{\large\bfseries\color{blue!70!black}}{\thesection.}{0.5em}{}
\titleformat{\subsection}{\normalsize\bfseries}{\thesubsection.}{0.5em}{}
\titleformat{\subsubsection}{\small\bfseries\color{gray!70!black}}{\thesubsubsection.}{0.5em}{}
\titlespacing*{\section}{0pt}{1.5ex}{1ex}
\titlespacing*{\subsection}{0pt}{1ex}{0.5ex}

\setlist[itemize]{nosep, leftmargin=1.5em}
\setlist[enumerate]{nosep, leftmargin=1.5em}

\tcbset{boxrule=0.5pt, left=4pt, right=4pt, top=4pt, bottom=4pt, fonttitle=\bfseries\small}

\newtcolorbox{definicja}[1][]{colback=blue!5, colframe=blue!50!black, title={#1}}
\newtcolorbox{uwagabox}[1][]{colback=orange!10, colframe=orange!70!black, title={#1}}
\newtcolorbox{wzorbox}[1][]{colback=gray!5, colframe=gray!50!black, title={#1}}
\newtcolorbox{przykladbox}[1][]{colback=purple!5, colframe=purple!50!black, title={#1}}
\usepackage{bookmark}
\IfFileExists{xurl.sty}{\usepackage{xurl}}{} % add URL line breaks if available
\urlstyle{same}
% fallback for those not using the hyperref driver hyperxmp:
\makeatletter
\@ifundefined{xmpquote}{\newcommand{\xmpquote}[1]{#1}}{}
\makeatother
\hypersetup{
  pdftitle={Handout W02 --- Język danych{,} pomiar i opis},
  pdfauthor={Marek Deja},
  pdflang={pl-PL},
  hidelinks,
  pdfcreator={LaTeX via pandoc}}

\title{Handout W02 --- Język danych, pomiar i opis}
\usepackage{etoolbox}
\makeatletter
\providecommand{\subtitle}[1]{% add subtitle to \maketitle
  \apptocmd{\@title}{\par {\large #1 \par}}{}{}
}
\makeatother
\subtitle{Mapa pojęć, wzorów i pytań}
\author{Marek Deja}
\date{2026/27}

\begin{document}
\maketitle

\section{1. Pytanie i główna
myśl}\label{pytanie-i-gux142uxf3wna-myux15bl}

\textbf{Pytanie:} Jak opisać samoocenę, obserwowany czas, powodzenie i
wybór kanału tak, aby miara odpowiadała skali i właściwemu mianownikowi?

Ta sama osoba może wysoko oceniać własne kompetencje i jednocześnie
długo wykonywać zadanie. Pomiar ankietowy oraz obserwacyjny są różnymi
oknami na zachowanie. Dobry opis nie zaciera tej różnicy i nie sprowadza
tabeli do jednej średniej.

\subsection{Tok rozumowania}\label{tok-rozumowania}

\begin{enumerate}
\def\labelenumi{\arabic{enumi}.}
\tightlist
\item
  Pytanie i zakres wniosku.\\
\item
  Jednostka, zmienne, skale i braki.\\
\item
  Estymata oraz jednostka.\\
\item
  Niepewność i dwa tory, jeśli dotyczą.\\
\item
  Efekt, ograniczenie i odpowiedź.
\end{enumerate}

\newpage

\section{2. Wzór i symbole}\label{wzuxf3r-i-symbole}

\[I_i=\frac{1}{m_i}\sum_{j=1}^{6}x^{*}_{ij}, \quad m_i\geq5; \qquad \bar{x}=\frac{1}{N}\sum x_i\]

{\def\LTcaptype{none} % do not increment counter
\begin{longtable}[]{@{}lll@{}}
\toprule\noalign{}
Symbol & Nazwa & Znaczenie \\
\midrule\noalign{}
\endhead
\bottomrule\noalign{}
\endlastfoot
\(x^{*}_{ij}\) & pozycja po ukierunkowaniu & punkty 1--5 \\
\(m_i\) & ważne pozycje osoby & co najmniej 5 \\
\(I_i\) & indeks samooceny & punkty 1--5 \\
\(\bar{x},Me\) & średnia i mediana & położenie rozkładu \\
\end{longtable}
}

Przeczytaj wzór zdaniem. Wskaż estymatę, punkt odniesienia, jednostkę i
składnik niepewności. Nie przypisuj p prawdopodobieństwa hipotezie ani
CI zakresu 95\% osób.

\newpage

\section{3. Dwa tory i efekt}\label{dwa-tory-i-efekt}

\textbf{Tor klasyczny.} Opis liczbowy obejmuje N, braki, miarę położenia
i zróżnicowania. Dla skośnego czasu zestawiamy średnią z medianą, SD z
IQR oraz histogram. Dla kategorii podajemy liczebność i odsetek z jawnym
mianownikiem.

\textbf{Tor permutacyjny/resampling.} Na etapie pomiaru nie używamy
testu do naprawiania złych wskaźników. Permutacja odpowie później na
pytanie o losowy związek przy zachowanych rozkładach. Najpierw musi być
jasne, co zostało zmierzone i które braki są dopuszczalne.

\textbf{Efekt i przedział.} Różnica między średnią i medianą opisuje
asymetrię, SD i IQR zróżnicowanie. Nie są to jeszcze efekty przyczynowe.
Jednostka miary pozwala ocenić, czy zmiana o dwie minuty lub pół punktu
jest ważna dla usługi.

\newpage

\section{4. Interpretacja i pytania}\label{interpretacja-i-pytania}

\textbf{Przykład.} Wielokrotny wybór kanałów może dać sumę odsetków
większą niż 100\%, ponieważ jedna osoba wybiera kilka opcji. Mianownik
powinien obejmować osoby z kompletnym zestawem odpowiedzi 0/1. Brak w
jednym kanale nie jest automatycznie odpowiedzią „nie''.

\textbf{Pułapki.} Pułapki obejmują kodowanie braku jako zera, odwrócenie
pozycji po policzeniu indeksu, średnią z nazw kategorii, wykres bez
jednostki, opis średniej jako wyniku większości i nazywanie kursowego
indeksu walidowanym testem.

\textbf{Pytanie do pracy.} Oceń cztery zdania z raportu: wskaż skalę,
mianownik, właściwą miarę oraz brakujące ograniczenie. Następnie
porównaj, co wnosi indeks samooceny, a co obserwowany czas i powodzenie.

\subsection{Dalszy krok}\label{dalszy-krok}

C03 zbuduje portret próby, a C04 przeprowadzi rekodację i indeks. W03
doda język niepewności, lecz zachowa rozróżnienie pomiaru od
wnioskowania.

\subsection{Literatura}\label{literatura}

\begin{itemize}
\tightlist
\item
  \href{https://publications.jrc.ec.europa.eu/repository/handle/JRC128415}{European
  Commission JRC (2022). DigComp 2.2}.
\item
  \href{https://www.ala.org/acrl/standards/ilframework}{Association of
  College and Research Libraries. Framework for Information Literacy}.
\item
  \href{https://www.webuse.org/pdf/Hargittai-BeyondLogsSurveys2002.pdf}{Hargittai,
  E. (2002). Beyond logs and surveys}.
\item
  \href{https://ris.utwente.nl/ws/portalfiles/portal/6972069/development.pdf}{van
  Deursen, Helsper \& Eynon. Internet Skills Scale}.
\end{itemize}

\end{document}
